\documentclass[
  12pt,
  openright,
  a4paper,
  twoside,
  english,
  brazil,
]{article}

\usepackage{mathtools}
\usepackage{amsfonts}
\usepackage{amssymb}
\usepackage{amsmath}
\usepackage[utf8]{inputenc}
\usepackage[portuguese]{babel}
\usepackage[T1]{fontenc}

\usepackage[top=0.5cm, bottom=2cm, left=2cm, right=2cm]{geometry}

\newcommand{\cqd}{\hfill $\square$}
\newenvironment{solution}[1][]{\noindent\textbf{Solução#1:} }{\cqd}
\newcounter{ex}
\newtheorem{exercise}[ex]{Exercício}

% -----------------------------------------------------------------------------

\title{Teoria dos Números: Lista 01}
\author{GIPMAT 2026 - Zoeh}
\date{}

\begin{document}

\maketitle

\begin{exercise}
  Prove por indução. \\
  Seja $a, b\in A$ e $m, n\in\mathbb N^*$. Então:

  \begin{enumerate}
    \item[(a)] $(a^m)^n=a^{m\cdot n}$
    \item[(b)] $(a\cdot b)^n=a^n \cdot b^n$
  \end{enumerate}
\end{exercise}

\begin{solution}
  a
\end{solution}

\begin{exercise}
  Mostre as seguintes fórmulas por indução

  \begin{enumerate}
    \item[(a)] $1^2+2^2+...+n^2=\frac{n(n+1)(2n+1)}{6}$
    \item[(b)] $\frac{1}{1\cdot 2}+\frac{1}{2\cdot
      3}+...+\frac{1}{n(n+1)}=\frac{n}{n+1}$
  \end{enumerate}
\end{exercise}

\begin{exercise}
  Uma progressão aritimética (PA) é uma sequência de números reais
  ($a_n$) tal que $a_1$ é dado e, para todo $n\in\mathbb N$, tem-se

  {\center$a_{n-1}=a_n+r$\endcenter}

  \hbox{onde $r$ é um número real fixo chamado razão}

  \begin{enumerate}
    \item[(a)] Mostre que $a_n=a_1+(n-1)r$
    \item[(b)] Se $a_n=a_1+a_2+...+a_n$. Mostre que
      $S_n=na_n+\frac{n(n-1)r}{2}=\frac{(a_1+a_n)n}{2}$
  \end{enumerate}
\end{exercise}

\begin{exercise}
  Prove por indução.

  \begin{enumerate}
    \item[(a)] $2^n>n$ para todo $n\in\mathbb N$
  \end{enumerate}
\end{exercise}

\end{document}
